\documentclass{beamer}

\usepackage{amsmath,amsfonts,amssymb}
\usepackage{physics}
\usepackage{bm}
\usepackage{quantikz}
\usepackage{listings}

\title{Quantum Phase Estimation}
\author{Morten Hjorth-Jensen}
\date{Additional notes spring 2026}

\lstset{
language=Python,
basicstyle=\ttfamily\small,
keywordstyle=\color{blue},
commentstyle=\color{green!60!black}
}

\begin{document}

\frame{\titlepage}

\section{Quantum Phase Estimation}

\begin{frame}{Eigenvalue Problem}

Suppose

\[
U|\psi\rangle=e^{2\pi i\phi}|\psi\rangle
\]

Goal

\[
\phi
\]

\end{frame}

%------------------------------------------------

\begin{frame}{Registers}

Two registers

\begin{itemize}
\item phase register
\item eigenstate register
\end{itemize}

\end{frame}

%------------------------------------------------

\begin{frame}{Initial State}

\[
|0\rangle^{\otimes t}|\psi\rangle
\]

\end{frame}

%------------------------------------------------

\begin{frame}{Superposition}

Apply Hadamards

\[
\frac{1}{2^{t/2}}
\sum_k |k\rangle |\psi\rangle
\]

\end{frame}

%------------------------------------------------

\begin{frame}{Controlled Powers}

Controlled

\[
U^{2^k}
\]

\end{frame}

%------------------------------------------------

\begin{frame}{State Before QFT}

\[
\frac{1}{2^{t/2}}
\sum_k e^{2\pi i k\phi}|k\rangle|\psi\rangle
\]

\end{frame}

%------------------------------------------------

\begin{frame}{Inverse QFT}

Inverse QFT converts phase to binary digits.

\end{frame}

%------------------------------------------------

\begin{frame}{Measurement}

Measurement yields

\[
\phi \approx 0.\phi_1\phi_2\dots
\]

\end{frame}

%================================================
\section{Error Analysis}

\begin{frame}{Finite Precision}

If

\[
\phi\neq k/2^t
\]

measurement yields nearest integer.

\end{frame}

%------------------------------------------------

\begin{frame}{Probability Distribution}

\[
P(y)=
\frac{1}{2^{2t}}
\left|
\sum_k e^{2\pi i k(\phi-y/2^t)}
\right|^2
\]

\end{frame}

%------------------------------------------------

\begin{frame}{Success Bound}

Probability

\[
P \ge 4/\pi^2
\]

\end{frame}

%================================================
\section{Iterative Phase Estimation}

\begin{frame}{Motivation}

Standard QPE requires many qubits.

Iterative version uses one ancilla.

\end{frame}

%------------------------------------------------

\begin{frame}{Iterative Circuit}

\begin{center}
\begin{quantikz}
\lstick{}&\gate{H}&\ctrl{1}&\meter{}\\
\lstick{psi}&\qw&\gate{U2k}&\qw
\end{quantikz}
\end{center}

\end{frame}

%================================================
\section{Hamiltonian Eigenvalue Estimation}

\begin{frame}{Hamiltonian Simulation}

If

\[
U=e^{-iHt}
\]

then

\[
U|\psi_k\rangle=e^{-iE_k t}|\psi_k\rangle
\]

\end{frame}

%------------------------------------------------

\begin{frame}{Energy Extraction}

\[
E_k =
\frac{2\pi\phi}{t}
\]

\end{frame}

%------------------------------------------------

\begin{frame}{Applications}

\begin{itemize}
\item quantum chemistry
\item materials science
\item nuclear physics
\end{itemize}

\end{frame}

%================================================
\section{Connection to VQE}

\begin{frame}{Comparison}

Phase estimation

\begin{itemize}
\item exact
\item deep circuits
\end{itemize}

VQE

\begin{itemize}
\item shallow circuits
\item classical optimization
\end{itemize}

\end{frame}

%================================================
\section{Python Simulation}

\begin{frame}[fragile]{QFT Matrix}

\begin{lstlisting}
import numpy as np

def qft_matrix(n):

    N=2**n
    omega=np.exp(2j*np.pi/N)

    F=np.zeros((N,N),dtype=complex)

    for j in range(N):
        for k in range(N):
            F[j,k]=omega**(j*k)

    return F/np.sqrt(N)
\end{lstlisting}

\end{frame}

%------------------------------------------------

\begin{frame}[fragile]{Example}

\begin{lstlisting}
n=3
F=qft_matrix(n)

state=np.zeros(2**n)
state[1]=1

result=F@state
print(result)
\end{lstlisting}

\end{frame}

%================================================
\section{Summary}

\begin{frame}{Summary}

\begin{itemize}
\item QFT efficiently computes Fourier transforms
\item central primitive in quantum algorithms
\item phase estimation extracts eigenvalues
\item widely used in quantum simulation
\end{itemize}

\end{frame}

\end{document}
